\section{Introduction}

Cloud platforms increasingly run distributed learning jobs on heterogeneous
and shared clusters.  Worker speeds, queueing delays, network variability, and
co-located interference create long-tail iteration times.  Tail-at-scale
analysis showed that a small fraction of slow components can dominate service
latency~\cite{dean2013tail}.  Cluster systems turned the same observation into
runtime mechanisms: speculative execution mitigates slow MapReduce tasks
\cite{zaharia2008late}, Mantri diagnoses and reacts to outliers in
MapReduce clusters~\cite{ananthanarayanan2010mantri}, and Dolly uses cloning to
reduce straggler impact~\cite{ananthanarayanan2013dolly}.  These systems
established a useful design principle: the scheduler should optimize the unit
that actually gates completion, not just the average work issued to the
cluster.  For synchronous distributed learning, that gate is revisited at every
iteration, so even modest worker-speed mismatch can repeatedly become a
barrier-latency problem.

Modern ML cluster schedulers apply this principle at the job and training-stack
layers.  Optimus models the marginal value of resources for distributed deep
learning jobs~\cite{peng2018optimus}.  Gandiva exploits introspection to
time-slice and pack deep learning workloads~\cite{xiao2018gandiva}.  Tiresias
schedules GPU jobs with learned job progress information~\cite{gu2019tiresias}.
Themis studies fair and efficient GPU sharing~\cite{mahajan2020themis}.  Gavel
makes heterogeneity an explicit cluster-scheduling dimension
\cite{narayanan2020gavel}.  Pollux co-adapts training batch size and cluster
allocation for goodput~\cite{qiao2021pollux}.  Sia extends goodput
optimization to heterogeneous ML clusters~\cite{subramanya2023sia}.  Recent
systems continue the same trajectory: RLTune combines learned prioritization
with solver-based allocation~\cite{dongare2025rltune}; Cuckoo optimizes
deadline-aware GPU job packing~\cite{zhang2025cuckoo}; Sailor automates
distributed training over dynamic heterogeneous clusters~\cite{strati2025sailor};
What-if analysis diagnoses training stragglers~\cite{lin2025whatifstragglers};
Holmes localizes irregularities in large GPU fleets~\cite{yao2025holmes};
Minder detects faulty machines during large-scale training~\cite{deng2025minder};
GREYHOUND targets fail-slow behavior in hybrid-parallel training
\cite{wu2025greyhound}; and Zeppelin balances variable-length LLM training
work~\cite{chen2026zeppelin}.  The common abstraction in these systems is that
heterogeneity is managed through job priority, resource allocation, packing,
placement, or diagnosis.  That abstraction is powerful, but it does not expose
the algebraic dependency among encoded rows inside a coded iteration.

Coded computing changes the completion condition of an iteration.  Gradient
coding uses algebraic redundancy so that a master can recover a full gradient
without waiting for every worker~\cite{tandon2017gradient}.  Short-Dot reduces
the cost of coded linear transforms through sparse coded products
\cite{dutta2016shortdot}.  Coded distributed learning generalizes this idea to
machine-learning iterations~\cite{lee2018speeding}.  Sparse coded matrix
multiplication further reduces coded-linear-algebra overhead
\cite{wang2018codedsparse}.  Approximate gradient coding studies heterogeneous
stragglers under relaxed recovery~\cite{song2025approxgc}.  Sparse flexible
coded learning (SFCL) creates the closest exact baseline for this paper: a
fixed sparse-flexible code matrix with multiple legal recovery
paths~\cite{zhang2025sparseflexible}.
These constructions move the barrier from ``all assigned work is complete'' to
``some algebraically sufficient row set is complete.''  This shift creates a
runtime question that code construction alone does not answer: among many legal
recovery paths, which one will become visible first under the current worker
speeds, dispatch delays, cancellation behavior, and sparse row costs.

The resulting scheduling problem is neither ordinary load balancing nor
general cluster placement.  In sparse flexible coded learning, the code defines
a row-span family rather than a single required worker set
\cite{zhang2025sparseflexible}.  An encoded row is useful only when it arrives
with other completed rows whose span contains the target gradient.  A
cost-balanced runtime can therefore place expensive sparse rows on fast workers
and still delay the first decodable row set, because row cost and row
decodability measure different things.  A generic heterogeneity-aware runtime
can also miss the same signal: cluster schedulers such as Gavel expose
heterogeneity as a job-to-resource placement problem~\cite{narayanan2020gavel},
whereas coded learning needs a row-prefix objective inside one iteration.  This
observation gives the background for our method.  The scheduler should preserve
the fixed
sparse-flexible code, predict how a candidate assignment changes the
first-decodable prefix, and enable adaptation only when the predicted prefix
reduction is large enough to pay for scheduling, extra-work, communication, and
cancellation overhead.

We present \textsc{G-port}, a guarded worker-service controller for sparse
flexible coded learning.  The name abbreviates Guarded Prefix-Oriented
Runtime: the controller keeps the code fixed and changes which row prefix is
exposed first.  Its coded arm assigns a target-specific importance score to
encoded rows, aggregates scores over row pairs, maps decode-critical pairs to
workers predicted to finish early, and runs the same row-span test as the
decoder to predict the first-decodable prefix.  A fixed online guard
enables the adaptive arm only when worker-speed mismatch, predicted
first-decode gain, prefix growth, and
rows-after-decode counters pass; otherwise \textsc{G-port} falls back to static
sparse-flexible placement.  The design keeps the
sparse-flexible code fixed and focuses on the worker-service runtime that
controls row exposure order.  Coded-only and alternative-fallback variants are
treated as ablations, not separate proposed systems.

The paper makes four contributions.

\begin{itemize}
  \item We identify first-decodable row exposure as a systems objective for
  sparse flexible coded learning.  The objective asks the runtime to expose an
  algebraically sufficient row prefix early, while leaving the code
  construction and decoder unchanged.
  \item We design \textsc{G-port}, a guarded worker-service controller with
  decode-aware row scores, finish-time estimates, first-prefix prediction, and
  static sparse-flexible fallback.  A fixed online guard enables the adaptive
  coded arm only when the predicted prefix reduction is worth its runtime cost.
  \item We specify the deployment semantics that make the design safe to use:
  reassignment changes row-to-worker order but not the recoverable row-span
  family, decoding is accepted only by the exact row-span test, and guard
  decisions are made before each iteration rather than after observing the
  outcome.
  \item We evaluate the design against scheduler-style baselines using
  trace-driven simulation and worker-service prototypes, including a direct
  three-node K3s deployment.  Across the main TCP, network-stressed TCP, and
  K3s worker-service comparisons, \textsc{G-port} remains positive at all tested
  worker counts; its bounded reissue path restores decode success under a
  single worker-close stress.
\end{itemize}
